\documentclass[11pt,a4paper]{article}
\usepackage[T1]{fontenc}
\usepackage[utf8]{inputenc}
\usepackage{lmodern}
\usepackage{amsmath,amssymb}
\usepackage{geometry}
\usepackage{microtype}
\geometry{margin=2.7cm}
\setlength{\parindent}{0pt}
\setlength{\parskip}{0.65em}

\begin{document}

\section*{Different Representations and Canonical Form of Simple Functions}

A simple function may admit different representations, corresponding to different choices of coefficients $a_k$ and measurable sets $A_k$.

For example, let $S=[0,1]$ and $\Sigma=\mathcal B([0,1])$. Then
\[
f(s)=\mathbf{1}_{[1/2,1]}(s)
=\mathbf{1}_{[1/2,2/3)}(s)
+\mathbf{1}_{[2/3,1]}(s).
\]

The equality follows because the two intervals on the right are disjoint and their union is $[1/2,1]$. Thus, the same function can be represented using different measurable sets.

A simple function is expressed in \emph{canonical form} if
\[
f(s)=\sum_{k=1}^{m}a_k\mathbf{1}_{A_k}(s),
\]
where
\[
a_1<a_2<\cdots<a_m,
\qquad
A_k=f^{-1}(\{a_k\}),
\quad k=1,\ldots,m.
\]

In this representation, the coefficients are precisely the distinct values attained by $f$, and the sets $A_k$ are their corresponding level sets. Consequently, these sets are pairwise disjoint and form a partition of $S$.

Every nonnegative simple function can be expressed in canonical form by collecting all points where it takes the same value. The corresponding level sets are measurable, pairwise disjoint, and form a finite partition of $S$. All values actually attained by $f$ must be included, including zero whenever it is attained.

\subsection*{Definition --- Integral of a Positive Simple Function}

We denote the collection of non-negative simple functions by $SF^+$. For $f\in SF^+$, represented as
\[
f=\sum_{k=1}^{m}a_k\mathbf{1}_{A_k},
\]
we define
\[
\mu_0(f):=\int_S f\,d\mu
:=\sum_{k=1}^{m}a_k\mu(A_k).
\]

This extends the definition of the integral of an indicator function, since for $f=\mathbf{1}_A$ we recover $\mu_0(f)=\mu(A)$.

\subsection*{Observation --- The Convention $0\cdot\infty=0$}

Let $S=\mathbb R$, $\Sigma=\mathcal B(\mathbb R)$, $\mu=\mathrm{Leb}$, and consider the zero function $f=0$.

We obtain
\[
\int_S f\,d\mu
=0\cdot\mu(\mathbb R)
=0\cdot\infty.
\]

Since the integral of the zero function must equal zero, we adopt the convention
\[
0\cdot\infty:=0.
\]

This ensures that zero coefficients contribute nothing to the integral, even when the corresponding sets have infinite measure.

Finally, since a simple function may admit different representations, we must verify that the value of $\mu_0(f)$ does not depend on the chosen representation.

This is the \emph{well-definedness of the integral of simple functions}, which will be established in the next homework.

\end{document}
